\hypertarget{listes}{%
\section{Listes}\label{listes}}

\hypertarget{une-liste-est-une-suxe9quence}{%
\subsection{Une liste est une
séquence}\label{une-liste-est-une-suxe9quence}}

Comme une chaîne de caractères, une \textbf{liste} est une séquence de
valeurs. Dans une chaîne, les valeurs sont des caractères ; dans une
liste, elles peuvent être de n\textquotesingle importe quel type. Les
valeurs qui composent une liste sont appelées des \textbf{éléments}
(\emph{items} ou \emph{elements}).

Il y a plusieurs façons de créer une nouvelle liste ; la plus simple
consiste à placer les éléments entre crochets (\texttt{{[}} et
\texttt{{]}}) :

\begin{Shaded}
\begin{Highlighting}[]
\NormalTok{devoirs }\OperatorTok{=}\NormalTok{ [}\StringTok{\textquotesingle{}maths\textquotesingle{}}\NormalTok{, }\StringTok{\textquotesingle{}programmation\textquotesingle{}}\NormalTok{, }\StringTok{\textquotesingle{}physique\textquotesingle{}}\NormalTok{]}
\NormalTok{nombres }\OperatorTok{=}\NormalTok{ [}\DecValTok{17}\NormalTok{, }\DecValTok{123}\NormalTok{]}
\NormalTok{vide }\OperatorTok{=}\NormalTok{ []}
\end{Highlighting}
\end{Shaded}

La première liste contient trois chaînes. La deuxième contient deux
entiers. La liste vide ne contient aucun élément.

Une liste à l\textquotesingle intérieur d\textquotesingle une autre
liste est dite \textbf{imbriquée} (\emph{nested}).

\begin{Shaded}
\begin{Highlighting}[]
\NormalTok{elements }\OperatorTok{=}\NormalTok{ [}\StringTok{\textquotesingle{}spam\textquotesingle{}}\NormalTok{, }\FloatTok{2.0}\NormalTok{, }\DecValTok{5}\NormalTok{, [}\DecValTok{10}\NormalTok{, }\DecValTok{20}\NormalTok{]]}
\end{Highlighting}
\end{Shaded}

\hypertarget{les-listes-sont-mutables}{%
\subsection{Les listes sont mutables}\label{les-listes-sont-mutables}}

La syntaxe pour accéder aux éléments d\textquotesingle une liste est la
même que pour les chaînes de caractères : l\textquotesingle opérateur de
crochet. L\textquotesingle indice indique quel élément vous voulez
(rappelez-vous que les indices commencent à 0) :

\begin{Shaded}
\begin{Highlighting}[]
\OperatorTok{\textgreater{}\textgreater{}\textgreater{}}\NormalTok{ devoirs[}\DecValTok{0}\NormalTok{]}
\CommentTok{\textquotesingle{}maths\textquotesingle{}}
\end{Highlighting}
\end{Shaded}

Contrairement aux chaînes de caractères, les listes sont
\textbf{mutables} (\emph{mutable}), car vous pouvez modifier
l\textquotesingle ordre des éléments d\textquotesingle une liste ou
réassigner un élément particulier. En utilisant
l\textquotesingle opérateur de crochet du côté gauche
d\textquotesingle une affectation, vous pouvez mettre à jour un élément
:

\begin{Shaded}
\begin{Highlighting}[]
\OperatorTok{\textgreater{}\textgreater{}\textgreater{}}\NormalTok{ nombres }\OperatorTok{=}\NormalTok{ [}\DecValTok{17}\NormalTok{, }\DecValTok{123}\NormalTok{]}
\OperatorTok{\textgreater{}\textgreater{}\textgreater{}}\NormalTok{ nombres[}\DecValTok{1}\NormalTok{] }\OperatorTok{=} \DecValTok{5}
\OperatorTok{\textgreater{}\textgreater{}\textgreater{}}\NormalTok{ nombres}
\NormalTok{[}\DecValTok{17}\NormalTok{, }\DecValTok{5}\NormalTok{]}
\end{Highlighting}
\end{Shaded}

L\textquotesingle élément de \texttt{nombres} qui se trouvait à
l\textquotesingle indice 1, qui était \texttt{123}, est maintenant
\texttt{5}.

Vous pouvez également penser à une liste comme à une relation entre les
indices et les éléments. Cette correspondance est appelée une
\textbf{séquence} ; c\textquotesingle est pourquoi les listes sont des
séquences.

L\textquotesingle opérateur \texttt{in} fonctionne également avec les
listes :

\begin{Shaded}
\begin{Highlighting}[]
\OperatorTok{\textgreater{}\textgreater{}\textgreater{}}\NormalTok{ devoirs }\OperatorTok{=}\NormalTok{ [}\StringTok{\textquotesingle{}maths\textquotesingle{}}\NormalTok{, }\StringTok{\textquotesingle{}programmation\textquotesingle{}}\NormalTok{, }\StringTok{\textquotesingle{}physique\textquotesingle{}}\NormalTok{]}
\OperatorTok{\textgreater{}\textgreater{}\textgreater{}} \StringTok{\textquotesingle{}maths\textquotesingle{}} \KeywordTok{in}\NormalTok{ devoirs}
\VariableTok{True}
\OperatorTok{\textgreater{}\textgreater{}\textgreater{}} \StringTok{\textquotesingle{}chimie\textquotesingle{}} \KeywordTok{in}\NormalTok{ devoirs}
\VariableTok{False}
\end{Highlighting}
\end{Shaded}

\hypertarget{parcours-dune-liste}{%
\subsection{Parcours d\textquotesingle une
liste}\label{parcours-dune-liste}}

La boucle \texttt{for} est un moyen élégant de parcourir les éléments
d\textquotesingle une liste. La syntaxe est la même que pour les chaînes
de caractères :

\begin{Shaded}
\begin{Highlighting}[]
\ControlFlowTok{for}\NormalTok{ devoir }\KeywordTok{in}\NormalTok{ devoirs:}
\NormalTok{    mettre\_a\_jour(devoir)}
\end{Highlighting}
\end{Shaded}

Cela fonctionne bien si vous avez seulement besoin de lire les éléments
de la liste. Mais si vous voulez écrire ou mettre à jour les éléments,
vous avez besoin des indices. Une façon courante de le faire est
d\textquotesingle utiliser les fonctions intégrées \texttt{range} et
\texttt{len} :

\begin{Shaded}
\begin{Highlighting}[]
\ControlFlowTok{for}\NormalTok{ i }\KeywordTok{in} \BuiltInTok{range}\NormalTok{(}\BuiltInTok{len}\NormalTok{(nombres)):}
\NormalTok{    nombres[i] }\OperatorTok{=}\NormalTok{ nombres[i] }\OperatorTok{*} \DecValTok{2}
\end{Highlighting}
\end{Shaded}

Cette boucle parcourt la liste et met à jour chaque élément.
\texttt{len} renvoie le nombre d\textquotesingle éléments de la liste.
\texttt{range} renvoie une liste d\textquotesingle indices de 0 à
\$n-1\$, où \$n\$ est la longueur de la liste. À chaque itération,
\texttt{i} prend l\textquotesingle indice de l\textquotesingle élément
suivant. L\textquotesingle instruction d\textquotesingle affectation
dans le corps utilise \texttt{i} pour lire l\textquotesingle ancien
élément et lui assigner le nouveau.

Une boucle \texttt{for} sur une liste vide ne s\textquotesingle exécute
jamais :

\begin{Shaded}
\begin{Highlighting}[]
\ControlFlowTok{for}\NormalTok{ x }\KeywordTok{in}\NormalTok{ []:}
    \BuiltInTok{print}\NormalTok{ Ceci ne s}\StringTok{\textquotesingle{}exécute jamais.}
\end{Highlighting}
\end{Shaded}

\hypertarget{opuxe9rations-sur-les-listes}{%
\subsection{Opérations sur les
listes}\label{opuxe9rations-sur-les-listes}}

L\textquotesingle opérateur \texttt{+} concatène les listes :

\begin{Shaded}
\begin{Highlighting}[]
\OperatorTok{\textgreater{}\textgreater{}\textgreater{}}\NormalTok{ a }\OperatorTok{=}\NormalTok{ [}\DecValTok{1}\NormalTok{, }\DecValTok{2}\NormalTok{, }\DecValTok{3}\NormalTok{]}
\OperatorTok{\textgreater{}\textgreater{}\textgreater{}}\NormalTok{ b }\OperatorTok{=}\NormalTok{ [}\DecValTok{4}\NormalTok{, }\DecValTok{5}\NormalTok{, }\DecValTok{6}\NormalTok{]}
\OperatorTok{\textgreater{}\textgreater{}\textgreater{}}\NormalTok{ c }\OperatorTok{=}\NormalTok{ a }\OperatorTok{+}\NormalTok{ b}
\OperatorTok{\textgreater{}\textgreater{}\textgreater{}}\NormalTok{ c}
\NormalTok{[}\DecValTok{1}\NormalTok{, }\DecValTok{2}\NormalTok{, }\DecValTok{3}\NormalTok{, }\DecValTok{4}\NormalTok{, }\DecValTok{5}\NormalTok{, }\DecValTok{6}\NormalTok{]}
\end{Highlighting}
\end{Shaded}

L\textquotesingle opérateur \texttt{*} répète une liste un nombre donné
de fois :

\begin{Shaded}
\begin{Highlighting}[]
\OperatorTok{\textgreater{}\textgreater{}\textgreater{}}\NormalTok{ [}\DecValTok{0}\NormalTok{] }\OperatorTok{*} \DecValTok{4}
\NormalTok{[}\DecValTok{0}\NormalTok{, }\DecValTok{0}\NormalTok{, }\DecValTok{0}\NormalTok{, }\DecValTok{0}\NormalTok{]}
\OperatorTok{\textgreater{}\textgreater{}\textgreater{}}\NormalTok{ [}\DecValTok{1}\NormalTok{, }\DecValTok{2}\NormalTok{, }\DecValTok{3}\NormalTok{] }\OperatorTok{*} \DecValTok{3}
\NormalTok{[}\DecValTok{1}\NormalTok{, }\DecValTok{2}\NormalTok{, }\DecValTok{3}\NormalTok{, }\DecValTok{1}\NormalTok{, }\DecValTok{2}\NormalTok{, }\DecValTok{3}\NormalTok{]}
\end{Highlighting}
\end{Shaded}

\hypertarget{tranches-de-listes}{%
\subsection{Tranches de listes}\label{tranches-de-listes}}

L\textquotesingle opérateur de tranche (\emph{slice}) fonctionne
également sur les listes :

\begin{Shaded}
\begin{Highlighting}[]
\OperatorTok{\textgreater{}\textgreater{}\textgreater{}}\NormalTok{ t }\OperatorTok{=}\NormalTok{ [}\StringTok{\textquotesingle{}a\textquotesingle{}}\NormalTok{, }\StringTok{\textquotesingle{}b\textquotesingle{}}\NormalTok{, }\StringTok{\textquotesingle{}c\textquotesingle{}}\NormalTok{, }\StringTok{\textquotesingle{}d\textquotesingle{}}\NormalTok{, }\StringTok{\textquotesingle{}e\textquotesingle{}}\NormalTok{, }\StringTok{\textquotesingle{}f\textquotesingle{}}\NormalTok{]}
\OperatorTok{\textgreater{}\textgreater{}\textgreater{}}\NormalTok{ t[}\DecValTok{1}\NormalTok{:}\DecValTok{3}\NormalTok{]}
\NormalTok{[}\StringTok{\textquotesingle{}b\textquotesingle{}}\NormalTok{, }\StringTok{\textquotesingle{}c\textquotesingle{}}\NormalTok{]}
\OperatorTok{\textgreater{}\textgreater{}\textgreater{}}\NormalTok{ t[:}\DecValTok{4}\NormalTok{]}
\NormalTok{[}\StringTok{\textquotesingle{}a\textquotesingle{}}\NormalTok{, }\StringTok{\textquotesingle{}b\textquotesingle{}}\NormalTok{, }\StringTok{\textquotesingle{}c\textquotesingle{}}\NormalTok{, }\StringTok{\textquotesingle{}d\textquotesingle{}}\NormalTok{]}
\OperatorTok{\textgreater{}\textgreater{}\textgreater{}}\NormalTok{ t[}\DecValTok{3}\NormalTok{:]}
\NormalTok{[}\StringTok{\textquotesingle{}d\textquotesingle{}}\NormalTok{, }\StringTok{\textquotesingle{}e\textquotesingle{}}\NormalTok{, }\StringTok{\textquotesingle{}f\textquotesingle{}}\NormalTok{]}
\end{Highlighting}
\end{Shaded}

Si vous omettez le premier indice, la tranche commence au début. Si vous
omettez le second, elle va jusqu\textquotesingle à la fin. Si vous
omettez les deux, la tranche est une copie de toute la liste :

\begin{Shaded}
\begin{Highlighting}[]
\OperatorTok{\textgreater{}\textgreater{}\textgreater{}}\NormalTok{ t[:]}
\NormalTok{[}\StringTok{\textquotesingle{}a\textquotesingle{}}\NormalTok{, }\StringTok{\textquotesingle{}b\textquotesingle{}}\NormalTok{, }\StringTok{\textquotesingle{}c\textquotesingle{}}\NormalTok{, }\StringTok{\textquotesingle{}d\textquotesingle{}}\NormalTok{, }\StringTok{\textquotesingle{}e\textquotesingle{}}\NormalTok{, }\StringTok{\textquotesingle{}f\textquotesingle{}}\NormalTok{]}
\end{Highlighting}
\end{Shaded}

Puisque les listes sont mutables, il est souvent utile de faire une
copie avant de les modifier.

Une tranche du côté gauche d\textquotesingle une affectation peut
modifier plusieurs éléments à la fois :

\begin{Shaded}
\begin{Highlighting}[]
\OperatorTok{\textgreater{}\textgreater{}\textgreater{}}\NormalTok{ t }\OperatorTok{=}\NormalTok{ [}\StringTok{\textquotesingle{}a\textquotesingle{}}\NormalTok{, }\StringTok{\textquotesingle{}b\textquotesingle{}}\NormalTok{, }\StringTok{\textquotesingle{}c\textquotesingle{}}\NormalTok{, }\StringTok{\textquotesingle{}d\textquotesingle{}}\NormalTok{, }\StringTok{\textquotesingle{}e\textquotesingle{}}\NormalTok{, }\StringTok{\textquotesingle{}f\textquotesingle{}}\NormalTok{]}
\OperatorTok{\textgreater{}\textgreater{}\textgreater{}}\NormalTok{ t[}\DecValTok{1}\NormalTok{:}\DecValTok{3}\NormalTok{] }\OperatorTok{=}\NormalTok{ [}\StringTok{\textquotesingle{}x\textquotesingle{}}\NormalTok{, }\StringTok{\textquotesingle{}y\textquotesingle{}}\NormalTok{]}
\OperatorTok{\textgreater{}\textgreater{}\textgreater{}}\NormalTok{ t}
\NormalTok{[}\StringTok{\textquotesingle{}a\textquotesingle{}}\NormalTok{, }\StringTok{\textquotesingle{}x\textquotesingle{}}\NormalTok{, }\StringTok{\textquotesingle{}y\textquotesingle{}}\NormalTok{, }\StringTok{\textquotesingle{}d\textquotesingle{}}\NormalTok{, }\StringTok{\textquotesingle{}e\textquotesingle{}}\NormalTok{, }\StringTok{\textquotesingle{}f\textquotesingle{}}\NormalTok{]}
\end{Highlighting}
\end{Shaded}

\hypertarget{muxe9thodes-de-listes}{%
\subsection{Méthodes de listes}\label{muxe9thodes-de-listes}}

Python fournit des méthodes qui s\textquotesingle appliquent aux listes.
Par exemple, \texttt{append} ajoute un nouvel élément à la fin
d\textquotesingle une liste :

\begin{Shaded}
\begin{Highlighting}[]
\OperatorTok{\textgreater{}\textgreater{}\textgreater{}}\NormalTok{ t }\OperatorTok{=}\NormalTok{ [}\StringTok{\textquotesingle{}a\textquotesingle{}}\NormalTok{, }\StringTok{\textquotesingle{}b\textquotesingle{}}\NormalTok{, }\StringTok{\textquotesingle{}c\textquotesingle{}}\NormalTok{]}
\OperatorTok{\textgreater{}\textgreater{}\textgreater{}}\NormalTok{ t.append(}\StringTok{\textquotesingle{}d\textquotesingle{}}\NormalTok{)}
\OperatorTok{\textgreater{}\textgreater{}\textgreater{}}\NormalTok{ t}
\NormalTok{[}\StringTok{\textquotesingle{}a\textquotesingle{}}\NormalTok{, }\StringTok{\textquotesingle{}b\textquotesingle{}}\NormalTok{, }\StringTok{\textquotesingle{}c\textquotesingle{}}\NormalTok{, }\StringTok{\textquotesingle{}d\textquotesingle{}}\NormalTok{]}
\end{Highlighting}
\end{Shaded}

\texttt{extend} prend une liste comme argument et ajoute tous ses
éléments :

\begin{Shaded}
\begin{Highlighting}[]
\OperatorTok{\textgreater{}\textgreater{}\textgreater{}}\NormalTok{ t1 }\OperatorTok{=}\NormalTok{ [}\StringTok{\textquotesingle{}a\textquotesingle{}}\NormalTok{, }\StringTok{\textquotesingle{}b\textquotesingle{}}\NormalTok{, }\StringTok{\textquotesingle{}c\textquotesingle{}}\NormalTok{]}
\OperatorTok{\textgreater{}\textgreater{}\textgreater{}}\NormalTok{ t2 }\OperatorTok{=}\NormalTok{ [}\StringTok{\textquotesingle{}d\textquotesingle{}}\NormalTok{, }\StringTok{\textquotesingle{}e\textquotesingle{}}\NormalTok{]}
\OperatorTok{\textgreater{}\textgreater{}\textgreater{}}\NormalTok{ t1.extend(t2)}
\OperatorTok{\textgreater{}\textgreater{}\textgreater{}}\NormalTok{ t1}
\NormalTok{[}\StringTok{\textquotesingle{}a\textquotesingle{}}\NormalTok{, }\StringTok{\textquotesingle{}b\textquotesingle{}}\NormalTok{, }\StringTok{\textquotesingle{}c\textquotesingle{}}\NormalTok{, }\StringTok{\textquotesingle{}d\textquotesingle{}}\NormalTok{, }\StringTok{\textquotesingle{}e\textquotesingle{}}\NormalTok{]}
\end{Highlighting}
\end{Shaded}

\texttt{sort} trie les éléments de la liste par ordre croissant :

\begin{Shaded}
\begin{Highlighting}[]
\OperatorTok{\textgreater{}\textgreater{}\textgreater{}}\NormalTok{ t }\OperatorTok{=}\NormalTok{ [}\StringTok{\textquotesingle{}d\textquotesingle{}}\NormalTok{, }\StringTok{\textquotesingle{}c\textquotesingle{}}\NormalTok{, }\StringTok{\textquotesingle{}e\textquotesingle{}}\NormalTok{, }\StringTok{\textquotesingle{}b\textquotesingle{}}\NormalTok{, }\StringTok{\textquotesingle{}a\textquotesingle{}}\NormalTok{]}
\OperatorTok{\textgreater{}\textgreater{}\textgreater{}}\NormalTok{ t.sort()}
\OperatorTok{\textgreater{}\textgreater{}\textgreater{}}\NormalTok{ t}
\NormalTok{[}\StringTok{\textquotesingle{}a\textquotesingle{}}\NormalTok{, }\StringTok{\textquotesingle{}b\textquotesingle{}}\NormalTok{, }\StringTok{\textquotesingle{}c\textquotesingle{}}\NormalTok{, }\StringTok{\textquotesingle{}d\textquotesingle{}}\NormalTok{, }\StringTok{\textquotesingle{}e\textquotesingle{}}\NormalTok{]}
\end{Highlighting}
\end{Shaded}

La plupart des méthodes de listes modifient l\textquotesingle argument
et retournent \texttt{None} ; elles sont conçues pour modifier des
listes en place et non pour créer de nouvelles listes.

\hypertarget{suppression-duxe9luxe9ments}{%
\subsection{Suppression
d\textquotesingle éléments}\label{suppression-duxe9luxe9ments}}

Il existe plusieurs façons de supprimer des éléments
d\textquotesingle une liste. Si vous connaissez l\textquotesingle indice
de l\textquotesingle élément que vous voulez supprimer, vous pouvez
utiliser \texttt{pop} :

\begin{Shaded}
\begin{Highlighting}[]
\OperatorTok{\textgreater{}\textgreater{}\textgreater{}}\NormalTok{ t }\OperatorTok{=}\NormalTok{ [}\StringTok{\textquotesingle{}a\textquotesingle{}}\NormalTok{, }\StringTok{\textquotesingle{}b\textquotesingle{}}\NormalTok{, }\StringTok{\textquotesingle{}c\textquotesingle{}}\NormalTok{]}
\OperatorTok{\textgreater{}\textgreater{}\textgreater{}}\NormalTok{ x }\OperatorTok{=}\NormalTok{ t.pop(}\DecValTok{1}\NormalTok{)}
\OperatorTok{\textgreater{}\textgreater{}\textgreater{}}\NormalTok{ t}
\NormalTok{[}\StringTok{\textquotesingle{}a\textquotesingle{}}\NormalTok{, }\StringTok{\textquotesingle{}c\textquotesingle{}}\NormalTok{]}
\OperatorTok{\textgreater{}\textgreater{}\textgreater{}}\NormalTok{ x}
\CommentTok{\textquotesingle{}b\textquotesingle{}}
\end{Highlighting}
\end{Shaded}

\texttt{pop\textasciigrave{}\ modifie\ la\ liste\ et\ renvoie\ l\textquotesingle{}élément\ qui\ a\ été\ supprimé.\ Si\ vous\ ne\ fournissez\ pas\ d\textquotesingle{}indice,\ il\ supprime\ et\ renvoie\ le\ dernier\ élément.\ \ Si\ vous\ ne\ connaissez\ pas\ l\textquotesingle{}indice\ mais\ que\ vous\ connaissez\ l\textquotesingle{}élément\ à\ supprimer,\ vous\ pouvez\ utiliser}remove\texttt{:\ \ ..\ code::\ python\ \ \ \ \ \ \textgreater{}\textgreater{}\textgreater{}\ t\ =\ {[}\textquotesingle{}a\textquotesingle{},\ \textquotesingle{}b\textquotesingle{},\ \textquotesingle{}c\textquotesingle{}{]}\ \ \ \ \ \textgreater{}\textgreater{}\textgreater{}\ t.remove(\textquotesingle{}b\textquotesingle{})\ \ \ \ \ \textgreater{}\textgreater{}\textgreater{}\ t\ \ \ \ \ {[}\textquotesingle{}a\textquotesingle{},\ \textquotesingle{}c\textquotesingle{}{]}\ \ La\ valeur\ de\ retour\ de}remove\texttt{est}None\texttt{.\ \ Pour\ supprimer\ plus\ d\textquotesingle{}un\ élément,\ vous\ pouvez\ utiliser\ une\ tranche\ avec\ une\ liste\ vide\ :\ \ ..\ code::\ python\ \ \ \ \ \ \textgreater{}\textgreater{}\textgreater{}\ t\ =\ {[}\textquotesingle{}a\textquotesingle{},\ \textquotesingle{}b\textquotesingle{},\ \textquotesingle{}c\textquotesingle{},\ \textquotesingle{}d\textquotesingle{},\ \textquotesingle{}e\textquotesingle{},\ \textquotesingle{}f\textquotesingle{}{]}\ \ \ \ \ \textgreater{}\textgreater{}\textgreater{}\ t{[}1:3{]}\ =\ {[}{]}\ \ \ \ \ \textgreater{}\textgreater{}\textgreater{}\ t\ \ \ \ \ {[}\textquotesingle{}a\textquotesingle{},\ \textquotesingle{}d\textquotesingle{},\ \textquotesingle{}e\textquotesingle{},\ \textquotesingle{}f\textquotesingle{}{]}\ \ Listes\ et\ chaînes\ -\/-\/-\/-\/-\/-\/-\/-\/-\/-\/-\/-\/-\/-\/-\/-\/-\ \ Une\ chaîne\ de\ caractères\ est\ une\ séquence\ de\ caractères\ et\ une\ liste\ est\ une\ séquence\ de\ valeurs,\ mais\ une\ liste\ de\ caractères\ n\textquotesingle{}est\ pas\ la\ même\ chose\ qu\textquotesingle{}une\ chaîne.\ Pour\ convertir\ une\ chaîne\ en\ une\ liste\ de\ caractères,\ vous\ pouvez\ utiliser}list\texttt{:\ \ ..\ code::\ python\ \ \ \ \ \ \textgreater{}\textgreater{}\textgreater{}\ s\ =\ "spam"\ \ \ \ \ \textgreater{}\textgreater{}\textgreater{}\ t\ =\ list(s)\ \ \ \ \ \textgreater{}\textgreater{}\textgreater{}\ t\ \ \ \ \ {[}\textquotesingle{}s\textquotesingle{},\ \textquotesingle{}p\textquotesingle{},\ \textquotesingle{}a\textquotesingle{},\ \textquotesingle{}m\textquotesingle{}{]}\ \ La\ fonction}list\texttt{sépare\ une\ chaîne\ en\ lettres\ individuelles.\ Si\ vous\ voulez\ diviser\ une\ chaîne\ en\ mots,\ vous\ pouvez\ utiliser\ la\ méthode}split\texttt{:\ \ ..\ code::\ python\ \ \ \ \ \ \textgreater{}\textgreater{}\textgreater{}\ s\ =\ "compter\ les\ mots\ de\ cette\ phrase"\ \ \ \ \ \textgreater{}\textgreater{}\textgreater{}\ t\ =\ s.split()\ \ \ \ \ \textgreater{}\textgreater{}\textgreater{}\ t\ \ \ \ \ {[}\textquotesingle{}compter\textquotesingle{},\ \textquotesingle{}les\textquotesingle{},\ \textquotesingle{}mots\textquotesingle{},\ \textquotesingle{}de\textquotesingle{},\ \textquotesingle{}cette\textquotesingle{},\ \textquotesingle{}phrase\textquotesingle{}{]}\ \ À\ l\textquotesingle{}inverse,}join\texttt{est\ l\textquotesingle{}opposé\ de}split\texttt{.\ Elle\ prend\ une\ liste\ de\ chaînes\ et\ concatène\ les\ éléments\ :\ \ ..\ code::\ python\ \ \ \ \ \ \textgreater{}\textgreater{}\textgreater{}\ t\ =\ {[}\textquotesingle{}compter\textquotesingle{},\ \textquotesingle{}les\textquotesingle{},\ \textquotesingle{}mots\textquotesingle{},\ \textquotesingle{}de\textquotesingle{},\ \textquotesingle{}cette\textquotesingle{},\ \textquotesingle{}phrase\textquotesingle{}{]}\ \ \ \ \ \textgreater{}\textgreater{}\textgreater{}\ delimiteur\ =\ \textquotesingle{}\ \textquotesingle{}\ \ \ \ \ \textgreater{}\textgreater{}\textgreater{}\ delimiteur.join(t)\ \ \ \ \ \textquotesingle{}compter\ les\ mots\ de\ cette\ phrase\textquotesingle{}\ \ Objets\ et\ valeurs\ -\/-\/-\/-\/-\/-\/-\/-\/-\/-\/-\/-\/-\/-\/-\/-\/-\ \ Si\ nous\ exécutons\ ces\ affectations\ :\ \ ..\ code::\ python\ \ \ \ \ \ a\ =\ "banane"\ \ \ \ \ b\ =\ "banane"\ \ Nous\ savons\ que}a\texttt{et}b\texttt{font\ référence\ à\ une\ chaîne\ de\ caractères,\ mais\ nous\ ne\ savons\ pas\ si\ elles\ font\ référence\ à\ la\ **même**\ chaîne.\ Il\ y\ a\ deux\ états\ possibles\ :\ \ D\textquotesingle{}un\ côté,}a\texttt{et}b\texttt{peuvent\ se\ référer\ à\ deux\ objets\ différents\ qui\ ont\ la\ même\ valeur.\ De\ l\textquotesingle{}autre,\ ils\ peuvent\ se\ référer\ au\ même\ objet.\ \ Pour\ vérifier\ si\ deux\ variables\ font\ référence\ au\ même\ objet,\ vous\ pouvez\ utiliser\ l\textquotesingle{}opérateur}is\texttt{:\ \ ..\ code::\ python\ \ \ \ \ \ \textgreater{}\textgreater{}\textgreater{}\ a\ =\ "banane"\ \ \ \ \ \textgreater{}\textgreater{}\textgreater{}\ b\ =\ "banane"\ \ \ \ \ \textgreater{}\textgreater{}\textgreater{}\ a\ is\ b\ \ \ \ \ True\ \ Dans\ cet\ exemple,\ Python\ ne\ crée\ qu\textquotesingle{}une\ seule\ chaîne\ de\ caractères,\ et}a\texttt{et}b\texttt{font\ référence\ au\ même\ objet.\ Mais\ lorsque\ vous\ créez\ deux\ listes,\ vous\ obtenez\ deux\ objets\ :\ \ ..\ code::\ python\ \ \ \ \ \ \textgreater{}\textgreater{}\textgreater{}\ a\ =\ {[}1,\ 2,\ 3{]}\ \ \ \ \ \textgreater{}\textgreater{}\textgreater{}\ b\ =\ {[}1,\ 2,\ 3{]}\ \ \ \ \ \textgreater{}\textgreater{}\textgreater{}\ a\ is\ b\ \ \ \ \ False\ \ Dans\ ce\ cas,\ on\ dit\ que\ les\ listes\ sont\ **équivalentes**\ (elles\ ont\ les\ mêmes\ valeurs),\ mais\ pas\ **identiques**\ (elles\ ne\ sont\ pas\ le\ même\ objet).\ \ Le\ crénelage\ (Aliasing)\ -\/-\/-\/-\/-\/-\/-\/-\/-\/-\/-\/-\/-\/-\/-\/-\/-\/-\/-\/-\/-\/-\/-\ \ Si}a\texttt{fait\ référence\ à\ un\ objet\ et\ que\ vous\ assignez}b
=
a\texttt{,\ alors\ les\ deux\ variables\ font\ référence\ au\ même\ objet\ :\ \ ..\ code::\ python\ \ \ \ \ \ \textgreater{}\textgreater{}\textgreater{}\ a\ =\ {[}1,\ 2,\ 3{]}\ \ \ \ \ \textgreater{}\textgreater{}\textgreater{}\ b\ =\ a\ \ \ \ \ \textgreater{}\textgreater{}\textgreater{}\ b\ is\ a\ \ \ \ \ True\ \ L\textquotesingle{}association\ d\textquotesingle{}une\ variable\ à\ un\ objet\ est\ appelée\ une\ **référence**.\ Si\ un\ objet\ a\ plus\ d\textquotesingle{}une\ référence,\ il\ a\ plus\ d\textquotesingle{}un\ nom,\ et\ on\ dit\ que\ l\textquotesingle{}objet\ est\ **crénelé**\ (*aliased*).\ \ Si\ l\textquotesingle{}objet\ crénelé\ est\ mutable,\ les\ modifications\ apportées\ à\ l\textquotesingle{}un\ des\ alias\ affectent\ l\textquotesingle{}autre\ :\ \ ..\ code::\ python\ \ \ \ \ \ \textgreater{}\textgreater{}\textgreater{}\ b{[}0{]}\ =\ 42\ \ \ \ \ \textgreater{}\textgreater{}\textgreater{}\ a\ \ \ \ \ {[}42,\ 2,\ 3{]}\ \ Bien\ que\ ce\ comportement\ puisse\ être\ utile,\ il\ est\ source\ d\textquotesingle{}erreurs.\ De\ manière\ générale,\ il\ est\ plus\ sûr\ d\textquotesingle{}éviter\ le\ crénelage\ lors\ de\ la\ manipulation\ de\ listes\ mutables.\ \ Arguments\ de\ listes\ -\/-\/-\/-\/-\/-\/-\/-\/-\/-\/-\/-\/-\/-\/-\/-\/-\/-\/-\ \ Lorsque\ vous\ passez\ une\ liste\ à\ une\ fonction,\ la\ fonction\ reçoit\ une\ référence\ à\ la\ liste.\ Si\ la\ fonction\ modifie\ un\ paramètre\ de\ la\ liste,\ la\ modification\ est\ visible\ par\ l\textquotesingle{}appelant.\ Par\ exemple,}supprimer\_premier\texttt{supprime\ le\ premier\ élément\ d\textquotesingle{}une\ liste\ :\ \ ..\ code::\ python\ \ \ \ \ \ def\ supprimer\_premier(t):\ \ \ \ \ \ \ \ \ del\ t{[}0{]}\ \ \ \ \ \ lettres\ =\ {[}\textquotesingle{}a\textquotesingle{},\ \textquotesingle{}b\textquotesingle{},\ \textquotesingle{}c\textquotesingle{}{]}\ \ \ \ \ supprimer\_premier(lettres)\ \ \ \ \ print(lettres)\ \ \ \ \ {[}\textquotesingle{}b\textquotesingle{},\ \textquotesingle{}c\textquotesingle{}{]}\ \ Débogage\ -\/-\/-\/-\/-\/-\/-\/-\ \ L\textquotesingle{}utilisation\ insouciante\ des\ listes\ (et\ d\textquotesingle{}autres\ objets\ mutables)\ peut\ entraîner\ de\ longues\ heures\ de\ débogage.\ Voici\ quelques\ pièges\ courants\ et\ des\ moyens\ de\ les\ éviter\ :\ \ 1.\ La\ plupart\ des\ méthodes\ de\ listes\ modifient\ l\textquotesingle{}argument\ et\ retournent}None\texttt{.\ C\textquotesingle{}est\ l\textquotesingle{}inverse\ des\ méthodes\ de\ chaînes,\ qui\ retournent\ une\ nouvelle\ chaîne\ et\ laissent\ l\textquotesingle{}originale\ intacte.\ 2.\ Choisissez\ un\ style\ de\ conception\ et\ tenez-s\textquotesingle{}y.\ Certaines\ fonctions\ modifient\ les\ objets\ en\ argument,\ d\textquotesingle{}autres\ créent\ de\ nouveaux\ objets.\ Essayez\ de\ combiner\ ces\ approches\ avec\ prudence.\ 3.\ Faites\ des\ copies\ à\ l\textquotesingle{}aide\ de\ tranches\ (}t{[}:{]}\texttt{)\ si\ vous\ souhaitez\ modifier\ une\ liste\ tout\ en\ gardant\ l\textquotesingle{}originale\ intacte.\ \ Glossaire\ -\/-\/-\/-\/-\/-\/-\/-\/-\ \ ..\ glossary::\ \ \ \ \ \ liste\ \ \ \ \ list\ \ \ \ \ \ \ \ \ Une\ séquence\ de\ valeurs.\ \ \ \ \ \ élément\ \ \ \ \ item\ \ \ \ \ \ \ \ \ L\textquotesingle{}une\ des\ valeurs\ d\textquotesingle{}une\ liste\ (ou\ d\textquotesingle{}une\ autre\ séquence).\ \ \ \ \ \ mutable\ \ \ \ \ mutable\ \ \ \ \ \ \ \ \ La\ propriété\ d\textquotesingle{}une\ séquence\ dont\ les\ éléments\ peuvent\ être\ modifiés.\ \ \ \ \ \ crénelage\ \ \ \ \ aliasing\ \ \ \ \ \ \ \ \ Le\ fait\ que\ deux\ ou\ plusieurs\ variables\ fassent\ référence\ au\ même\ objet.\ \ \ \ \ \ délimiteur\ \ \ \ \ delimiter\ \ \ \ \ \ \ \ \ Un\ caractère\ ou\ une\ chaîne\ utilisée\ pour\ indiquer\ où\ une\ chaîne\ doit\ être\ divisée.\ \ \ Exercices\ -\/-\/-\/-\/-\/-\/-\/-\/-\ \ ..\ topic::\ Exercice\ 1\ \ \ \ \ \ Écrivez\ une\ fonction\ appelée}somme\_cumulee\texttt{qui\ prend\ une\ liste\ de\ nombres\ et\ renvoie\ la\ somme\ cumulée\ ;\ c\textquotesingle{}est-à-dire\ une\ nouvelle\ liste\ où\ le\ \$i\$-ème\ élément\ est\ la\ somme\ des\ premiers\ \$i+1\$\ éléments\ de\ la\ liste\ originale.\ Par\ exemple,\ la\ somme\ cumulée\ de}{[}1,
2, 3{]}\texttt{est}{[}1, 3,
6{]}\texttt{.\ \ ..\ topic::\ Exercice\ 2\ \ \ \ \ \ Écrivez\ une\ fonction\ nommée}centre\texttt{qui\ prend\ une\ liste\ et\ renvoie\ une\ nouvelle\ liste\ sans\ le\ premier\ et\ le\ dernier\ éléments.\ Par\ exemple,\ si\ vous\ passez}{[}1,
2, 3, 4{]}\texttt{,\ elle\ doit\ renvoyer}{[}2,
3{]}\texttt{.\ \ ..\ topic::\ Exercice\ 3\ \ \ \ \ \ Écrivez\ une\ fonction\ appelée}est\_trie\texttt{qui\ prend\ une\ liste\ en\ paramètre\ et\ renvoie}True\texttt{si\ la\ liste\ est\ triée\ par\ ordre\ croissant,\ et}False\texttt{sinon.\ Par\ exemple,}est\_trie({[}1,
2,
2{]})\texttt{devrait\ renvoyer}True\texttt{et}est\_trie({[}\textquotesingle b\textquotesingle,
\textquotesingle a\textquotesingle{]})\texttt{devrait\ renvoyer}False\texttt{.\ \ ..\ topic::\ Exercice\ 4\ \ \ \ \ \ Écrivez\ une\ fonction\ appelée}est\_anagramme\texttt{qui\ prend\ deux\ chaînes\ de\ caractères\ et\ renvoie}True\texttt{si\ elles\ sont\ des\ anagrammes\ (si\ elles\ contiennent\ les\ mêmes\ lettres\ avec\ les\ mêmes\ fréquences).\ \ ..\ topic::\ Exercice\ 5\ \ \ \ \ \ Écrivez\ une\ fonction\ appelée}contient\_doublons\texttt{qui\ prend\ une\ liste\ et\ renvoie}True``
si il y a un élément qui apparaît plus d\textquotesingle une fois. Elle
ne doit pas modifier la liste originale.

\textbf{Exercice 6}

Cet exercice est basé sur le fameux « Paradoxe des anniversaires ».
Écrivez une fonction qui génère 23 dates aléatoires (représentées par
des entiers entre 1 et 365) et vérifie s\textquotesingle il y a des
doublons. En exécutant cela des milliers de fois, estimez la probabilité
qu\textquotesingle il y ait au moins deux personnes nées le même jour
dans un groupe de 23 personnes.
